\documentclass[UTF8,zihao=-4]{ctexart}

\usepackage[a4paper,top=2.4cm,bottom=2.4cm,left=2.25cm,right=2.25cm]{geometry}
\usepackage{amsmath,amssymb}
\usepackage{booktabs,longtable,array,tabularx}
\usepackage{enumitem}
\usepackage{xcolor}
\usepackage{hyperref}
\usepackage{fancyhdr}

\definecolor{inkblue}{RGB}{34,76,112}
\definecolor{softgray}{RGB}{246,248,250}
\hypersetup{
  colorlinks=true,
  linkcolor=inkblue,
  urlcolor=inkblue,
  citecolor=inkblue,
  pdfauthor={数学建模学习小组},
  pdftitle={数学建模阶段学习笔记与论文写作模板}
}

\pagestyle{fancy}
\fancyhf{}
\fancyhead[L]{数学建模阶段学习笔记与论文写作模板}
\fancyhead[R]{\thepage}
\setlength{\headheight}{15pt}
\setlength{\parindent}{2em}
\setlength{\parskip}{0.25em}
\setlist{nosep,leftmargin=2.2em}

\title{\bfseries 数学建模阶段学习笔记与论文写作模板}
\author{队长、组员1、组员2}
\date{2026年8月}

\begin{document}

\maketitle

\begin{abstract}
本报告记录本队现阶段的算法学习、论文精读和写作模板。第一部分依照思维导图，
按三人分工整理十三个学习单元，列出基本原理、求解步骤、代码入口和适用边界；
第二部分选取2024年A题一等奖论文A053《基于目标优化的板凳龙形态调节》，分析其
摘要、假设、建模、求解、结果解释与附录的写法；最后给出本队以后写论文时采用的
章节模板。文中的“已学习”指已能复述原理、说明步骤并列出代码逻辑，不等同于已经
完成全部算法的限时独立复现。

\par\noindent\textbf{关键词：}数学建模；学习笔记；板凳龙；论文结构；写作模板
\end{abstract}

\tableofcontents
\newpage

